% ============================================================
\part{Netzwerk und Kommunikation}
% ============================================================

\chapter{Wir wollen zwei Programme miteinander sprechen lassen}

\kapitelziel{
Du verstehst, wie Linux-Prozesse miteinander kommunizieren können
und kennst die grundlegenden IPC-Mechanismen.
}
\kapitelstruktur

\section{Situation}
Dein Sensor-Programm liest Werte, dein Navigations-Programm
braucht diese Werte -- aber sie laufen in getrennten Prozessen.
Wie kommen die Daten von A nach B?

\begin{aufgabeBox}
Lass zwei Python-Programme über eine Pipe miteinander
kommunizieren.
\end{aufgabeBox}

\begin{problemBox}
Prozesse haben getrennte Speicherbereiche. Sie können nicht
einfach Variablen teilen. Kommunikation braucht explizite Kanäle.
\end{problemBox}

\begin{wissenBox}
Wir brauchen: Pipes, Named Pipes (FIFOs), Sockets und den
Unterschied zwischen synchroner und asynchroner Kommunikation.
\end{wissenBox}

\section{Pipes -- der einfachste Weg}

\begin{lstlisting}[style=bash]
# Unix-Pipe: Ausgabe von Programm A als Eingabe von Programm B
python3 sensor_ausgabe.py | python3 navigation_eingabe.py

# Named Pipe (FIFO) -- persistenter Kanal
mkfifo /tmp/sensor_kanal

# Terminal 1: Daten senden
python3 -c "
import time
with open('/tmp/sensor_kanal', 'w') as f:
    for i in range(10):
        f.write(f'{i * 0.1:.2f}\n')
        f.flush()
        time.sleep(0.5)
"

# Terminal 2: Daten empfangen
python3 -c "
with open('/tmp/sensor_kanal', 'r') as f:
    for line in f:
        print(f'Empfangen: {line.strip()}')
"
\end{lstlisting}

\section{Sockets -- netzwerkfähige Kommunikation}

\begin{lstlisting}[style=python]
# sensor_server.py -- Sensorwerte senden
import socket, time, random

server = socket.socket(socket.AF_INET, socket.SOCK_STREAM)
server.bind(('localhost', 9000))
server.listen(1)
print("Warte auf Verbindung...")
conn, addr = server.accept()

while True:
    wert = round(random.uniform(0.1, 2.0), 2)
    conn.send(f"{wert}\n".encode())
    time.sleep(0.5)
\end{lstlisting}

\begin{lstlisting}[style=python]
# navigations_client.py -- Sensorwerte empfangen
import socket

client = socket.socket(socket.AF_INET, socket.SOCK_STREAM)
client.connect(('localhost', 9000))

while True:
    daten = client.recv(1024).decode().strip()
    abstand = float(daten)
    if abstand < 0.25:
        print(f"WARNUNG: Hindernis bei {abstand} m")
    else:
        print(f"Frei: {abstand} m")
\end{lstlisting}

\begin{loesungBox}
Für einfache Kommunikation auf demselben Rechner: Named Pipe.
Für netzwerkfähige Kommunikation (Roboter ↔ Laptop): Socket.
Für komplexe Systeme mit vielen Teilnehmern: ROS2 Topics.
\end{loesungBox}

\begin{merksatzBox}
IPC (Inter-Process Communication) ist das Fundament vernetzter
Systeme. Pipes sind einfach, aber eingeschränkt. Sockets sind
flexibel und netzwerkfähig. ROS2 baut auf DDS auf -- einem
industriellen Nachrichtenprotokoll, das Sockets weit übertrifft.
\end{merksatzBox}

\begin{robotikBox}
ROS2's DDS-Middleware ist ein hochentwickeltes IPC-System:
Topics sind Named Pipes, die automatisch entdeckt werden
(\textit{Discovery}), mit Typsicherheit und Quality-of-Service-Garantien.
Der Schritt von Socket zu ROS2-Topic ist konzeptionell klein.
\end{robotikBox}

% ── Kapitel 22 ────────────────────────────────────────────────
\chapter{Wir wollen Nachrichten statt direkter Befehle senden}

\kapitelziel{
Du verstehst das Publish-Subscribe-Muster und kannst es
in Python implementieren -- als Vorbereitung für ROS2-Topics.
}
\kapitelstruktur

\section{Situation}
Sensor, Navigation und Anzeige sollen alle Abstandswerte bekommen.
Mit direkten Verbindungen (A→B, A→C) wird der Code zu einem
Geflecht von Abhängigkeiten.

\begin{aufgabeBox}
Implementiere ein einfaches Publish-Subscribe-System, bei dem
ein Sender Daten veröffentlicht und mehrere Empfänger diese
unabhängig voneinander erhalten.
\end{aufgabeBox}

\begin{analogie}
Publish-Subscribe ist wie ein Radiosender: Der Sender (Publisher)
sendet ins Leere. Wer zuhören will (Subscriber), schaltet ein.
Der Sender kümmert sich nicht darum, wie viele zuhören.
\end{analogie}

\begin{lstlisting}[style=python]
# pubsub.py -- Minimales Publish-Subscribe in Python
import threading, time, random
from collections import defaultdict

class MessageBus:
    """Einfacher Message-Bus: Publish-Subscribe."""
    def __init__(self):
        self._subscribers = defaultdict(list)

    def subscribe(self, topic, callback):
        self._subscribers[topic].append(callback)

    def publish(self, topic, message):
        for cb in self._subscribers[topic]:
            cb(message)

# Bus-Instanz (geteilt zwischen allen)
bus = MessageBus()

# Subscriber-Callbacks
def navigation_callback(msg):
    abstand = msg['abstand']
    if abstand < 0.25:
        print(f"  Navigation: STOPP (Abstand: {abstand} m)")

def anzeige_callback(msg):
    print(f"  Anzeige: Abstand = {msg['abstand']:.2f} m")

# Registrieren
bus.subscribe('/sensor/abstand', navigation_callback)
bus.subscribe('/sensor/abstand', anzeige_callback)

# Publisher (Sensor) in eigenem Thread
def sensor_publisher():
    while True:
        wert = round(random.uniform(0.1, 2.0), 2)
        bus.publish('/sensor/abstand', {'abstand': wert, 'zeit': time.time()})
        time.sleep(1.0)

t = threading.Thread(target=sensor_publisher, daemon=True)
t.start()
time.sleep(5)
\end{lstlisting}

\begin{merksatzBox}
Publisher kennen ihre Subscriber nicht. Subscriber kennen ihren
Publisher nicht. Beide kennen nur den Topic-Namen. Diese
\textit{lose Kopplung} macht das System flexibel -- neue Subscriber
können jederzeit hinzukommen, ohne Publisher zu ändern.
\end{merksatzBox}

\begin{robotikBox}
Dieser Code ist strukturell identisch mit ROS2:
\begin{lstlisting}[style=python]
# ROS2-Equivalent:
self.publisher = self.create_publisher(Float32, '/sensor/abstand', 10)
self.subscription = self.create_subscription(
    Float32, '/sensor/abstand', self.navigation_callback, 10)
\end{lstlisting}
Das Konzept -- Publish-Subscribe über benannte Topics --
ist exakt dasselbe. ROS2 fügt Netzwerkfähigkeit,
Typsicherheit und Discovery hinzu.
\end{robotikBox}
